\subsection{Orientation, one formula, and the bridge to mechanics}

The equation
\[
r=\frac{p}{1+e\cos\theta}
\]
places the nearest point on the positive polar axis. Rotating the entire conic by
an angle $\theta_0$ gives
\[
\boxed{r=\frac{p}{1+e\cos(\theta-\theta_0)}}.
\]
The parameters have separate geometric roles:
\[
\begin{array}{c|c}
p & \text{overall focal scale},\\
e & \text{shape and conic type},\\
\theta_0 & \text{orientation in the plane}.
\end{array}
\]

\begin{center}
\begin{tikzpicture}[scale=0.8,>=stealth]
    \draw[axis,->] (-5.2,0)--(5.5,0) node[right] {$x$};
    \draw[axis,->] (0,-4.1)--(0,4.3) node[above] {$y$};
    \begin{scope}[rotate=32]
        \draw[subset] (-0.7,0) ellipse (3.7 and 2.0);
        \fill[geometricSubsetColor] (0,0) circle (2.2pt) node[below left] {focus};
        \draw[blue!70!black,line width=1.1pt,->] (0,0)--(3.0,0)
            node[midway,above] {periapsis direction};
    \end{scope}
    \draw[orange!85!black,line width=1pt,->] (1.2,0)
        arc[start angle=0,end angle=32,radius=1.2];
    \node at (1.35,0.55) {$\theta_0$};
\end{tikzpicture}
\end{center}

The same formula now separates the three qualitative regimes:
\[
\begin{array}{c|c|c}
e & \text{conic} & \text{focal motion suggested by the geometry}\\
\hline
0\leq e<1 & \text{circle or ellipse} & \text{bounded},\\
e=1 & \text{parabola} & \text{threshold between bounded and unbounded},\\
e>1 & \text{hyperbola} & \text{unbounded passage}.
\end{array}
\]

\begin{corebox}[What has been achieved]
The chapter did not merely list four equations. It changed the viewpoint until
three apparently different curves became outcomes of one radial law. This is the
analytic-geometric preparation needed before a later mechanics course derives
the same law from a central force.
\end{corebox}

\begin{summarybox}[Choose coordinates adapted to the question]
Cartesian coordinates are best for centres, axes, vertices, and complete
symmetric pictures. Polar coordinates centred at a focus are best for radial
distance, orientation, eccentricity, and orbital interpretation. Analytic
geometry works by choosing the representation in which the geometry speaks most
clearly.
\end{summarybox}
